\chapter{Introduction}\label{introduction}

After the inception of blockchain as a distributed ledger in 2009 with
the invention of Bitcoin, the concept of \textbf{blockchain
interoperability (BI)} soon followed and started to receive attention
\citep{Belchior_2021}. It is the ability of different blockchain
networks to communicate and exchange data, or currency and assets
without needing a centralised arbiter
\citep{kishore2025interoperability}. This is motivated by the
realisation that blockchains have become isolated silos or systems
incapable of communicating with their outside environment.
Interoperability permits technology to develop beyond its limitation to
scalability, reduces fragmentation in the ecosystem, and increases the
use of liquidity in different platforms \citep{belenkov2025review}
\citep{Saleh2024}. Without this, innovation will remain limited as
digital assets and information remain stagnant and isolated within their
own blockchain networks.

The early protocols for BI started development within 2015-2016. They
are mainly focused on the atomic exchanage of digital asset or tokens
between different networks
\citep{wilson2026exploringblockchaininteroperabilityframeworks}. These
protocols, such as sidechains, notary schemes, and hashed time-lock
contracts (HTLC), are categorised as \textbf{asset-based} since their
original goal is to facilitate the exchange of value and cryptocurrency
through an autonomous and synchronised way \citep{wang2021sok}.

When the blockchain usage moved beyond cryptocurrency to complicated
applications like provenance, tracking, and Internet of Things (IoT),
new protocols aimed to handle arbitrary data evolved \citep{Kotey_2024}.
As data from many applications quickly grow in volume, the need to
prioritise sharing arbitrary information and state transitions between
networks has become more evident \citep{wang2021sok}
\citep{kishore2025interoperability}. Since the earlier blockchains are
treated as isolated silos, the mere exchange of value is insufficient
for sectors like supply chain and healthcare that need wider
coordination and sharing of information
\citep{wilson2026exploringblockchaininteroperabilityframeworks}
\citep{madine2021appxchain}. This need led to the creation of
cross-blockchain smart contracts and messaging frameworks allowing
decentralised applications (dApps) to coordinate and verify information
from different chain without undergoing a centralised arbiter
\citep{Belchior_2021}.

Recognising the limitations of single-purpose protocols, developers
created unified interoperability protocols capable of handling both
assets and arbitrary data. The \textbf{Interblockchain (IBC)} protocol
is the leading example of this. It handles both digital assets and
arbitraty data by means of its payload-agnostic design; information is
delivered as byte-encoded packets whatever it contains
\citep{cosmos2024ibcintro} \citep{goes2020interblockchain}. This works
by separating the transport layer (TAO) (that focuses on transportation,
authentication, and arrangment of data) and the application layer (that
gives meaning to each payload) \citep{ibc2025howworks}.

For assets, IBC uses guidelines like (InterChain Standards) ICS-20 for
locking the original asset and minting the corresponding vouchers in the
target blockchain \citep{quasar2023litepaper}
\citep{ibc2024transportapp} \citep{goes2020interblockchain}. Meanwhile
for arbitrary data and application logic, ICS-27 (Interchain Accounts)
and ICS-31 (Interchain Queries) are used in order to execute commands or
read states from other network without needing a centralised arbiter
\citep{quasar2023litepaper} \citep{ibc2024transportapp}. IBC serves as a
general-purpose communication layer similar to TCP/IP of the Internet
that allows heterogenous blockchains to exchange whatever type of
information in a way that is safe and trust-minimised.

\section{Project Context}\label{project-context}

The \textbf{Inter-Blockchain Communication (IBC) Protocol} is considered
as the most used and widely adopted interoperability protocol that can
handle both digital assets and arbitrary data
\citep{ibc2024transportapp} \citep{Belchior_2021}. As a general-purpose
message-passing protocol, it serves as the connective tissue of more
than 100 different blockchains (where 107 chains uses reference
implementation of IBC-Go) to cooperate in a way that is secure,
permissionless, and trust-minimised. Its payload-agnostic design enables
it to deliver whatever information that is encoded in bytes --- from
simple token transfers using ICS-20 standard to more complex arbitrary
data and application logic like ICS-27 (Interchain Accounts) and ICS-31
(Interchain Queries). Because of its wide adoption and ability to serve
as a universal standard like TCP/IP of Internet, it is the leading
solution in creating an Internet of Blockchains that supports different
workflows in the whole ecosystem \citep{kishore2025interoperability}.

However, despite IBC\textquotesingle s widespread adoption, a
fundamental question remains unanswered: Do IBC\textquotesingle s
asset-transfer and data-transfer behaviours share the same core
semantics (unified), are assets semantics merely built on top of data
transport without fundamental integration (staked), or do these
behaviours interfere with each other in ways that create observable
anomalies (conflicting)?

\textbf{Semantic interoperability}, or the ability of different
blockchain networks to understand exchanged data and use it meaningfully
in their own systems, depends on answering this question
\citep{Saleh2024} \citep{Belchior_2021}. Without a clear understanding
of how asset and data behaviours relate within a unified protocol like
IBC, developers cannot be certain that cross-chain operations will
behave predictably. At the same time, auditors cannot guarantee that
security properties hold across both asset transfers and arbitrary
message passing.

\textbf{Formal modelling} answers the problem of semantic
interoperability by providing a mathematical and clear description of
the characteristics of the systems that removes ambiguity or
misconceptions on how data should be understood and processed
\citep{verma2022formalmethods}. It serves as the foundation in ensuring
that application-specific semantics are maintained amidst the
differences in architecture and consensus mechanisms \citep{Saleh2024}
\citep{Belchior_2021}. Formal verification like model checking, provides
irrefutable evidence to the developers and auditors that the
coordination of the behaviours of assets and arbitrary data remain
predictable and safe within a unified protocol like the IBC.

Currently, formalisation on IBC is focused on its TAO using tools like
TLA+ in order to verify the correctness, safety, and liveness of its
protocol \citep{Wei_2023}. There is a need to formalise its semantic
interoperability layer because although ICS dictates guidelines for
applications, there is insufficient mathematical foundation that
verifies the coordination of the behaviours of assset and arbitrary data
within system.

This research will use \textbf{Process Algebra} to model the behaviour
of IBC protocol.

\section{Purpose and Description}\label{purpose-and-description}

The main purpose of this research is to formally specify and analyse the
observable behaviours of the IBC protocol using process algebra and
labelled transition systems (LTS). Specifically, this research will:

\begin{enumerate}
\tightlist
\item
  Produce the first complete process algebra specification of
  IBC\textquotesingle s core protocols (ICS-004 for data transport and
  ICS-020 for asset transfer).
\item
  Extract LTS behavioural models from these specifications.
\item
  Analyse whether IBC\textquotesingle s asset-transfer and data-transfer
  behaviours are unified, stacked, or conflicting.
\end{enumerate}

Unlike the prior formal analyses of IBC using TLA+ trace-based
verification, this research uses process algebra which is a formalism
uniquely suited to capturing concurrent interaction, branching
structure, and silent abstraction. This provides a complementary
perspective that can reveal behavioural anomalies that trace-based
methods might miss.

Aside from developers and auditors, this will give dApp developers
assurance and predictability in developing ccross-chain applications.
For chain operator and validator, this will reduce development risk and
ensures stability of the system against dynamic runtime errors like
deadlock or livelock \citep{Kotey_2024} \citep{Duan2018}. Lastly,
end-users can ensure the integrity and security of their asset and
transaction in a trust-minimised environment \citep{Kotey_2024}.

To guide this research, the following questions are posed:

\begin{enumerate}
\tightlist
\item
  How can IBC (including both asset transfer and data transfer) be
  formally specified using process algebra?
\item
  What obervable behaviours (traces, ordering constraints, branching
  structure, deadlocks, synchronisation patterns) emerge from the LTS
  model of IBC?
\item
  Within the IBC specification, are the behavioural patterns for asset
  transfer and data transfer: a. Unified (they share the same core
  semantics)? b. Stacked (asset semantics are built on top of data
  transport without fundamental integration?) c. Conflicting (they
  interfere with each other in ways that cause observable behavioural
  anomalies?)
\end{enumerate}

\section{Objectives}\label{objectives}

To achive the main purpose of this research, the following specific
objectives must be accomplished:

\begin{enumerate}
\tightlist
\item
  Identify and define the atomic actions, alphabets, communication
  ports, and data types necessary to model IBC\textquotesingle s ICS-004
  and ICS-020 in process algebra.
\item
  Formally specify IBC\textquotesingle s core interaction patterns using
  process algebra operators.
\item
  Extract LTS models from the process algebra specifications, capturing
  all reachable states and observable transitions.
\item
  Analyse the LTS models to identify traces, ordering constraints,
  branching structures, deadlock states, and synchrnisation pattern for
  both asset-transfer and data-transfer scenarios.
\item
  Compare the behavioural patterns for asset transfer and data transfer
  within IBC specification to determine whether they are unified,
  stacked, or conflicting.
\item
  Document the process algebra specification and LTS models in a
  reusable form suitable for future comparative studies with other
  interoperability protocols like Polkadot XCM and Axelar.
\end{enumerate}

\section{Scope and Limitations}\label{scope-and-limitations}

The formal model abstracts away low-level details such as cryptographic
verification (Merkle proofs), light client synchronisation, and relayer
behaviour

Economic incentives (gas costs, slashing, liquidity constraints) are not
modelled

The analysis focuses on observable behaviour only; security properties
beyond what is observable (e.g., double-spending resistance) are not
verified

Only IBC\textquotesingle s core ICS-004 and ICS-020 standards are
modelled; other application standards (ICS-27, ICS-31) are excluded

No empirical validation or actual protocol instance execution is
performed
